
\subsection{Invariant Primitives and \texttt{InvCell} Verification} \label{sec:invcell}

Just as Rust's standard library implements \verb`Cell` via \verb`UnsafeCell`,
in \verus we can implement and verify \verb`InvCell` via our 
\verb`UnsafeCell`-equivalent, \verb`PCell`.
To do this, though, we first need to introduce our \emph{invariant primitives}.

To see what these are for, consider what happens when we
try to implement \verb`InvCell<T>` using \verb`UnsafeCell<T>`.
From the API, we know that we need to be able to write
even when we only have access to a shared reference \verb`&InvCell<T>`,
but writing to the underlying \verb`UnsafeCell<T>` requires exclusive ownership
of the \verb`PermData<T>` object.

Once again, we run into this problem of trying to gain exclusive ownership of something
that is shared. However, we have pushed the problem one layer down---to the \emph{ghost}
layer, and this is where \verus introduces its trusted invariant primitives to escape
the problem.

The two primitives are called \verb`LocalInvariant<G>` and
\verb`AtomicInvariant<G>`.\footnote{Both these types also have additional type parameters used to specify the invariant as a boolean predicate on \texttt{G}.}
Each one allows the user to store a (ghost) object \verb`G`; each one allows the user to perform
a ghost operation called \emph{opening the invariant}, where they obtain temporary, exclusive
ownership over the \verb`G`. For example, this snippet shows how the implementation of
\verb`InvCell<T>::replace` temporarily gains access to the \verb`PermData<T>` object:

\begin{lstlisting}[language=Verus,style=VerusLineNos]
impl InvCell<T> {
    pub fn replace(&self, val: T) -> T {
        requires(self.wf() && self.inv(val));
        ensures(|old_val| self.inv(old_val));

        let r;

        // Opens the invariant `&self.perm_inv` which has type `LocalInvariant<PermData<T>>`.
        // Opening the invariant is a ghost operation, and it binds to the ghost variable `perm`.
        open_local_invariant!(&self.perm_inv => perm => {
            // The code inside, however, is executable. This is where we actually perform
            // the write, using ownership of the ghost `perm` object, of type `PermData<T>`.
            r = self.pcell.replace(&mut perm, val);
        });

        r // Return the old value.
    }
}
\end{lstlisting}

The difference between the two primitives is that \verb`LocalInvariant<G>` is restricted for use
on a single thread: it does not implement \verb`Send` or \verb`Sync`, the traits Rust uses to mark thread-safety.
\verb`AtomicInvariant<G>` is thread-safe, and it does implement these traits: however,
this comes with an additional requirement, that the invariant may only be opened
for atomic operations.
Since \verb`InvCell` (like Rust's \verb`Cell`) is for single-threaded use,
we use \verb`LocalInvariant<V>` here. 

The reader might wonder what happens if we attempt to nest calls to the invariant-opening
operation, \verb`open_local_invariant`. It would certainly be unsound if we could open the same
\verb`LocalInvariant<G>` object twice, and obtain double-ownership of the `T`.
Indeed, \verus generates extra verification conditions to disallow such things
by tracking which invariants are ``open'' at a given time. These verification conditions
are designed to be lightweight, and they have no impact on our SMT generation
for cases outside of those which use the low-level invariant APIs.
